% Tratándose de una aplicación de detección de posturas para realizar ejercicio físico,
% se debería tratar el tema del ejercicio físico y su importancia en la salud.
% También se puede mencionar la gamificación y su uso en aplicaciones para fomentar el ejercicio.
% Otra cosa que se puede mencionar es la importancia de la ergonomía en el ejercicio físico,
% y cómo la detección de posturas puede ayudar a mejorar la ergonomía en el ejercicio
% y prevenir lesiones y hacer más fácil el ejercicio.

\chapter{Antecedentes}
\label{cap:antecedentes}
En este capítulo se presentan los antecedentes relevantes para el desarrollo de este proyecto,
entendiendo como antecedentes aquellos estudio e investigaciones previas que ayude a contextualizar
el tema del trabajo de fin de grado. También se presentan de forma breve los aspectos teóricos previos
del problema.

Tratándose de una aplicación de detección de posturas para realizar ejercicio físico, se tratará el tema
del ejercicio físico y su importancia en la salud, así como de la gamificación como herramienta para fomentar
el ejercicio.

\section{Ejercicio físico}
\label{sec:ejercicio-fisico}
% Qué es
% Beneficios
% Carecemos de ejercicio y sedentarismo
% Consecuencias del sedentarismo
Se entiende por ejercicio físico a \enquote{cualquier actividad física que mejora y mantiene la
aptitud física, la salud y el bienestar del individuo} \parencite{WikipediaEjercicioFisico}. El
ejercicio físico es una parte fundamental de un estilo de vida saludable y tiene numerosos beneficios
para la salud física y mental. Entre algunos de estos beneficios se encuentran la reducción del estrés,
ansiedad y depresión \parencites{Guterman2023Jul}{marquez1995beneficios}, el aumento de la calidad de
vida y el grado de independencia de personas mayores \parencite{Morel2009EjercicioAdultoMayor} o los
efectos protectores contra muchas enfermedades neurodegenerativas y neuromusculares
\parencite{grondard2005regular}.

Sin embargo, la falta de ejercicio físico es un problema creciente en la sociedad actual. Tener un
estilo de vida sedentario, caracterizado por la falta de actividad física regular, puede tener
consecuencias graves para la salud. El sedentario se ha relacionado con un mayor riesgo de
enfermedades crónicas como la obesidad, la diabetes tipo 2, enfermedades cardiovasculares, ciertos
tipos de cáncer, cansancio, estrés y ansiedad, entre otros \parencite{GarcíaMatamoros_2019}.

% https://www.ine.es/ss/Satellite?param1=PYSDetalleFichaIndicador&c=INESeccion_C&param3=1259937499084&p=1254735110672&pagename=ProductosYServicios%2FPYSLayout&cid=1259944495973&L=0
Según una encuesta realizada por el Instituto Nacional de Estadística (INE) en el año 2022, solo el
\textbf{37,2\%} de la población española (mayores de 15 años) realiza ejercicio físico de forma
regular. Del resto de la población, el \textbf{27,4\%} dice llevar un estilo de vida sedentario.
Aunque este porcentaje ha disminuido en los últimos años, sigue siendo un problema de salud pública
importante. Es por esto que fomentar el ejercicio físico y reducir el sedentarismo es una prioridad
para mejorar la salud y el bienestar de la población.

\diagrama{
    nombre = {Recomendaciones de ejercicio físico},
    fuente = {Ministerio de Sanidad del Gobierno de España},
    archivo = {memoria/fotos/Recomendaciones_ActividadFisica_poblacion_adulta.pdf},
    etiqueta = {fig:recomendaciones-ejercicio-fisico},
    width = 1\textwidth,
}


\section{Ludificación}
\label{sec:ludificacion}
% Qué es
% Cómo se puede utilzar en aplicaciones de ejercicio
% Ejemplo de aplicación de ejercicio con gamificación
% Decir que el tfg cae dentro de la gamificación
La ludificación (más conocida por el anglicismo \textit{gamificación}) es el uso de elementos y
técnicas de diseño de juegos en contextos no lúdicos para motivar y mejorar la participación de
los usuarios. Su objetivo es hacer que ciertas actividades sean más atractivas, placenteras y
motivadoras, utilizando elementos como puntos, niveles, recompensas, desafíos y fomentando la
competitividad entre usuarios.

La gamificación se ha utilizado en diversas áreas, como la educación, el marketing y la salud,
para mejorar la participación y el compromiso de los usuarios. En el ámbito del ejercicio físico,
la gamificación puede ser una herramienta efectiva para motivar a las personas a realizar actividad
física de forma regular. Existen varios ejemplos de aplicaciones de ejercicio físico que
utilizan la gamificación para fomentar la actividad física, una de las más recientes siendo
\textit{Ring Fit Adventure}, un videojuego de Nintendo que combina el ejercicio físico con
elementos de juego, como desafíos, recompensas y niveles \footnote{\url{https://ringfitadventure.nintendo.com/}}.
Después de buscar información sobre la aplicación y sus resultados, se ha comprobado que a un gran
grupo de personas les ha motivado a realizar ejercicio físico de forma regular durante un largo
periodo de tiempo, lo que demuestra la efectividad de la gamificación en este ámbito%
\footnote{\url{https://cgicoffee.com/blog/2021/07/a-full-year-of-fitness-with-ring-fit-adventure}}%
\footnote{\url{https://www.ign.com/articles/how-ring-fit-adventure-transformed-lives-those-who-beat}}.

Esto es algo que se pretende conseguir con el proyecto de fin de grado, utilizando la ludificación
como una herramienta para motivar a los usuarios a realizar ejercicio físico de forma regular.
